% Append-only section fragment. Paste below the preceding section in the master document.
% This fragment intentionally contains no preamble and no document boundaries.

\section{Programme Identity and Public Catalogue Evidence}
\label{sec:programme-identity-catalogue}

Pilot Fabric 0.41 introduces a provider-independent programme-identity layer for
production screen perception. The purpose of this layer is to identify a likely
programme from permitted screen evidence without making the stronger and often
incorrect claim that the programme is currently playing, available in the user's
region, included in a subscription, or accessible by the household.

\subsection{Operating Sequence}

The programme-identity pipeline applies the following ordered process:

\begin{enumerate}
  \item An owner starts the visible Observer or submits a television photograph.
  \item Apple Vision performs local OCR and classification on the mother device.
  \item The source image is deleted immediately after local analysis.
  \item The screen-fusion layer accepts title candidates only on a recognised
        media surface such as Netflix, YouTube, or another media application.
  \item Obvious controls, times, prices, URLs, profile prompts, and other interface
        text are removed from the candidate set.
  \item The multi-frame perception timeline requires the same OCR title to occur
        in at least two distinct recent frames before treating it as stable.
  \item A stable text-only candidate may be submitted to the bounded public
        programme catalogue.
  \item Pilot accepts the catalogue identity only when exactly one result has the
        same normalized title. Ambiguous and approximate matches fail closed.
  \item The phone controller displays the title, provenance, attribution, and the
        separate catalogue, provider, and playback verification states.
\end{enumerate}

This ordering prevents a large piece of interface copy, a subtitle, an advert,
or a single OCR error from becoming a silently asserted programme identity.

\subsection{Evidence Classes and Trust Boundaries}

Pilot does not collapse all metadata into one confidence score. It records four
different identity sources:

\begin{description}
  \item[Signed provider adapter] A locally approved, persona-bound adapter has
        signed normalized provider and playback facts. This is the only source
        that can verify current playback.
  \item[Official provider metadata] An official provider API has verified a
        provider content record. This establishes metadata, not playback or
        entitlement.
  \item[Public programme catalogue] One strict exact title match has established
        a public catalogue record. It does not establish the active application,
        playback, availability, subscription, or household access.
  \item[Repeated OCR] The same title has appeared in at least two distinct local
        observations. This is stable screen evidence but remains provider-unverified.
\end{description}

The resulting programme record exposes separate Boolean fields for catalogue
identity, provider verification, current playback, availability, and account
entitlement. A true value in one field never promotes another field implicitly.

\subsection{Public Catalogue Adapter}

The initial public adapter uses the TVmaze search interface. Pilot sends only a
bounded title string after local repetition has established it. Unicode accents
and punctuation are normalized for exact comparison, but fuzzy similarity is not
accepted as identity. If two catalogue entries share the same normalized title,
Pilot reports an ambiguous match and leaves the title unverified.

Successful and unsuccessful results are cached on the mother brain for twenty-four
hours. The cache is limited to one hundred entries and is stored inside the
owner-only runtime area. Network responses are limited to 500 kilobytes and a
ten-second timeout. The retained record includes TVmaze attribution and a link to
the public catalogue source. TVmaze data is used under its published CC BY-SA
terms; the production product must preserve attribution and review share-alike
obligations for any redistributed derivative dataset.

\subsection{Privacy and Data Minimisation}

No source photograph, video frame, raw audio, provider token, account identifier,
profile name, viewing history, or credential is sent to the public catalogue.
Only the repeated programme-title candidate is transmitted. Source pixels remain
local and are deleted after analysis. The structured perception timeline retains
at most twelve hashed frames, and a context transition prevents title evidence
from leaking from one application or input into another.

The public lookup is therefore an optional evidence-acquisition step rather than
a cloud-vision dependency. When the Internet or catalogue is unavailable, repeated
OCR remains useful and is clearly labelled as local observation rather than
verified catalogue identity.

\subsection{Runtime Integration}

The implementation is divided into reusable components:

\begin{itemize}
  \item \texttt{ScreenUnderstandingFusion} extracts and ranks bounded OCR title
        candidates on media surfaces.
  \item \texttt{MultiFrameScreenPerception} requires repeated agreement, separates
        evidence classes, and resets context after a verified surface transition.
  \item \texttt{ProviderMetadataResolver} resolves official provider identifiers
        and performs the strict public-catalogue lookup when repeated evidence exists.
  \item \texttt{TrustedProviderTelemetry} independently verifies signed playback
        and entitlement facts from explicitly approved adapters.
  \item The private companion displays identity provenance, catalogue status,
        playback status, attribution, uncertainty, and retained subtitle history.
\end{itemize}

The live capture path first considers explicit device playback metadata and the
previous stable repeated-OCR result. It then resolves programme evidence and
queries signed telemetry using the application-provider hint, ensuring that the
generic catalogue adapter cannot accidentally replace Netflix, YouTube, or another
provider identity in the trusted telemetry lookup.

\subsection{Failure Behaviour}

Pilot fails closed when the title is seen once, OCR confidence is too low, the
surface is not a media surface, the response is too large, the catalogue is
unavailable, no exact title exists, or several exact records exist. In each case
the observation can remain visible as uncertain local evidence, but no provider,
playback, subscription, or entitlement claim is created.

\subsection{Verification}

Automated verification covers official YouTube enrichment, provider identifiers,
strict public-catalogue exact matching, ambiguous-title rejection, twenty-four-hour
cache reuse, media-only OCR candidate extraction, interface-text rejection,
multi-frame OCR stability, evidence-source priority, context-transition isolation,
and the rule that a catalogue match cannot verify current playback. Pilot Fabric
0.41 passes 310 automated Fabric tests.

\subsection{Remaining Production Work}

The reusable evidence contract is complete, but production coverage remains open.
Pilot still requires authorized signed-in provider adapters, film and episode-level
catalogues, programme-guide schedules, advertisement and live-event identity,
permitted caption interfaces, and field qualification across television models,
regions, languages, and application layouts. Those adapters must use the same
separation between identity, playback, availability, and entitlement.
